% GENERATED FILE. Edit canonical lessons, then run npm run generate:books.
% canonical-source-hash: fnv1a64:43239db719c6e313
% canonical-lessons: TE-C222-sarada, TE-C222-adhunika, TE-C222-spastamaina, TE-C222-bhinnamaina, TE-C222-nyayamaina

\chapter{Fun, Modern, Clear, Different, Fair}
\label{ch:te-fun-modern-clear-different-fair}
\hlchaptermodality{222}

\begin{chapteropening}
\textbf{By the end of this chapter:} \emph{I can say fun, modern, clear, different and fair.}

Complete the last lesson of chapter 222: I can say fun, modern, clear, different and fair.
\end{chapteropening}

\section[\texorpdfstring{\te{సరదా} (saradā)}{సరదా (saradā)}]{\te{సరదా} (saradā) — fun; a fondness}

\label{lesson:TE-C222-sarada}



\begin{quote}
Before the new one: say the Telugu for round, then the Telugu for sharp.
\end{quote}

\subsection*{You'll want to know: \te{సరదా}}
\textbf{\te{సరదా}} — \emph{saradā} — \textquotedblleft{}fun; a fondness\textquotedblright{}.

\begin{grammarlens}[title={What you've built}]
One more word for this chapter.
\end{grammarlens}

\subsection*{Your turn}
\begin{itemize}
\raggedright
  \item \emph{Say it:} \emph{saradā}
  \item \emph{Say it:} \emph{saradā}, once more
  \item \emph{Say it:} say \emph{saradā}
  \item \emph{Recall:} say the Telugu for round, then the Telugu for sharp, then say \emph{saradā} again
\end{itemize}

\begin{tcolorbox}[breakable,colback=teal!4,colframe=teal!35!black,title={Before you move on}]
What does \te{సరదా} mean? (\textquotedblleft{}Fun; a fondness\textquotedblright{}.) Say it once more.
\end{tcolorbox}
\section[\texorpdfstring{\te{ఆధునిక} (ādhunika)}{ఆధునిక (ādhunika)}]{\te{ఆధునిక} (ādhunika) — modern}

\label{lesson:TE-C222-adhunika}



\begin{quote}
Before the new one: say the Telugu for sharp, then the Telugu for fun.
\end{quote}

\subsection*{You'll want to know: \te{ఆధునిక}}
\textbf{\te{ఆధునిక}} — \emph{ādhunika} — \textquotedblleft{}modern\textquotedblright{}.

\begin{grammarlens}[title={What you've built}]
One more word for this chapter.
\end{grammarlens}

\subsection*{Your turn}
\begin{itemize}
\raggedright
  \item \emph{Say it:} \emph{ādhunika}
  \item \emph{Say it:} \emph{ādhunika}, once more
  \item \emph{Say it:} say \emph{ādhunika}
  \item \emph{Recall:} say the Telugu for sharp, then the Telugu for fun, then say \emph{ādhunika} again
\end{itemize}

\begin{tcolorbox}[breakable,colback=teal!4,colframe=teal!35!black,title={Before you move on}]
What does \te{ఆధునిక} mean? (\textquotedblleft{}Modern\textquotedblright{}.) Say it once more.
\end{tcolorbox}
\section[\texorpdfstring{\te{స్పష్టమైన} (spaṣṭamaina)}{స్పష్టమైన (spaṣṭamaina)}]{\te{స్పష్టమైన} (spaṣṭamaina) — clear}

\label{lesson:TE-C222-spastamaina}



\begin{quote}
Before the new one: say the Telugu for fun, then the Telugu for modern.
\end{quote}

\subsection*{You'll want to know: \te{స్పష్టమైన}}
\textbf{\te{స్పష్టమైన}} — \emph{spaṣṭamaina} — \textquotedblleft{}clear\textquotedblright{}.

\begin{grammarlens}[title={What you've built}]
One more word for this chapter.
\end{grammarlens}

\subsection*{Your turn}
\begin{itemize}
\raggedright
  \item \emph{Say it:} \emph{spaṣṭamaina}
  \item \emph{Say it:} \emph{spaṣṭamaina}, once more
  \item \emph{Say it:} say \emph{spaṣṭamaina}
  \item \emph{Recall:} say the Telugu for fun, then the Telugu for modern, then say \emph{spaṣṭamaina} again
\end{itemize}

\begin{tcolorbox}[breakable,colback=teal!4,colframe=teal!35!black,title={Before you move on}]
What does \te{స్పష్టమైన} mean? (\textquotedblleft{}Clear\textquotedblright{}.) Say it once more.
\end{tcolorbox}
\section[\texorpdfstring{\te{భిన్నమైన} (bhinnamaina)}{భిన్నమైన (bhinnamaina)}]{\te{భిన్నమైన} (bhinnamaina) — different}

\label{lesson:TE-C222-bhinnamaina}



\begin{quote}
Before the new one: say the Telugu for modern, then the Telugu for clear.
\end{quote}

\subsection*{You'll want to know: \te{భిన్నమైన}}
\textbf{\te{భిన్నమైన}} — \emph{bhinnamaina} — \textquotedblleft{}different\textquotedblright{}.

\begin{grammarlens}[title={What you've built}]
One more word for this chapter.
\end{grammarlens}

\subsection*{Your turn}
\begin{itemize}
\raggedright
  \item \emph{Say it:} \emph{bhinnamaina}
  \item \emph{Say it:} \emph{bhinnamaina}, once more
  \item \emph{Say it:} say \emph{bhinnamaina}
  \item \emph{Recall:} say the Telugu for modern, then the Telugu for clear, then say \emph{bhinnamaina} again
\end{itemize}

\begin{tcolorbox}[breakable,colback=teal!4,colframe=teal!35!black,title={Before you move on}]
What does \te{భిన్నమైన} mean? (\textquotedblleft{}Different\textquotedblright{}.) Say it once more.
\end{tcolorbox}
\section[\texorpdfstring{\te{న్యాయమైన} (nyāyamaina)}{న్యాయమైన (nyāyamaina)}]{\te{న్యాయమైన} (nyāyamaina) — fair, righteous}

\label{lesson:TE-C222-nyayamaina}



\begin{quote}
Before the new one: say the Telugu for clear, then the Telugu for different.
\end{quote}

\subsection*{You'll want to know: \te{న్యాయమైన}}
\textbf{\te{న్యాయమైన}} — \emph{nyāyamaina} — \textquotedblleft{}fair, righteous\textquotedblright{}.

\begin{grammarlens}[title={What you've built}]
One more word for this chapter.
\end{grammarlens}

\subsection*{Your turn}
\begin{itemize}
\raggedright
  \item \emph{Say it:} \emph{nyāyamaina}
  \item \emph{Say it:} \emph{nyāyamaina}, once more
  \item \emph{Say it:} say \emph{nyāyamaina}
  \item \emph{Recall:} say the Telugu for clear, then the Telugu for different, then say \emph{nyāyamaina} again
\end{itemize}

\begin{tcolorbox}[breakable,colback=teal!4,colframe=teal!35!black,title={Before you move on}]
What does \te{న్యాయమైన} mean? (\textquotedblleft{}Fair, righteous\textquotedblright{}.) Say it once more.
\end{tcolorbox}
